\begin{table*}[t]\centering\footnotesize
\caption{Research questions, systematization findings and evidence. Experimental follow-ups test the comparison rules; the SoK contribution is the record-level synthesis.}\label{tab:rq-findings}
\begin{tabularx}{\textwidth}{@{}p{.20\textwidth}YY@{}}\toprule
Question & Finding & Evidence in the main paper\\\midrule
RQ1: targets and observables & A common signal can serve confidentiality, attribution, evasion or authorized diagnostics; polarity changes the protected interest. & Complete core-family map (\S\ref{sec:corpus-map}, Table~\ref{tab:core-family-map}).\\
RQ2: origin versus access & B-size cover does not bind A-call order; plaintext re-emission creates a fresh B trace; shared A-call projections do not establish unchanged state. These distinctions identify the evidence required to reuse a result. & NetEcho/SplitDecoy, WRA/AHB and dispatcher read/write records (\S\ref{sec:discussion}); Exps~1--2.\\
RQ3: portability & A task holdout does not establish topology, hardware or model transfer; mixed configurations can preserve shortcuts. & Generator-specific holdouts (\S\ref{sec:evaluation}); customer-grouped analysis (\S\ref{sec:r5}).\\
RQ4: protection, state and cost & Bind inference, utility and cost to one scoped operating point. The 20-family audit gives R8 Bound 3/17 for component-or-task utility and 2/16 under a task-only sensitivity. & Twenty-family R1--R9 matrix; Exps~3--4 leakage/state/cost contrasts (\S\ref{sec:controlled-placement}). The same audit locates the utility boundary of our demonstrations.\\\bottomrule
\end{tabularx}\end{table*}
